\chapter{Programming code}
\label{app:code}

All analyses were carried out in Python. The programs are submitted as separate files in a ZIP archive. They are listed below in the order in which they are run; each reads the outputs of the previous ones from a common working folder. All random numbers are generated from fixed seeds. The FBref and Transfermarkt data are collected by the programs and are not redistributed.

\begin{table}[ht]
	\centering
	\footnotesize
	\caption{Programming files}
	\label{tab:code}
	\begin{tabular}{p{4.6cm}p{10.2cm}}
		\toprule
		File & Purpose \\
		\midrule
		\texttt{test\_fbref\_collection\_v2.py} & Registers the Primeira Liga in \texttt{soccerdata}, collects four FBref tables (standard, shooting, playing time, miscellaneous) for 2017/18 to 2024/25 and merges them. \\
		\texttt{collect\_tm\_v2.py} & Collects Transfermarkt squad lists for each club and season, links FBref players to Transfermarkt players within each season, downloads the market value histories and extracts the end-of-season values. \\
		\texttt{build\_panel.py} & Removes rows that cannot be matched (same-name players at one club), merges the rows of players with several clubs in a season, applies the sample filters, builds the per-90 statistics, shrunk rates and position indicators, and scales the variables. \\
		\texttt{make\_analysis\_panel.py} & Attaches market values, resolves players linked to two spellings, defines one identity per player and constructs the outcome for consecutive-season pairs. \\
		\texttt{crosswalk\_audit.py} & Lists the weakest FBref--Transfermarkt links for manual review. \\
		\texttt{fcmd\_mdnc.py} & Module with the clustering model: type-specific dissimilarities, the fuzzy $C$-medoids algorithm with noise cluster, fixed or estimated weights, the Xie--Beni index and the selection routines. \\
		\texttt{cluster\_one\_season.py} & Development check of the clustering model on a single season. \\
		\texttt{align\_and\_select.py} & Clusters every season, aligns the archetypes across seasons with the Hungarian algorithm on the distance between their profiles, builds the membership transitions and the diagnostics used to choose $K$. \\
		\texttt{pairs\_lib.py} & Module that builds the table of season pairs and the baselines, and the tests used by the regression programs. \\
		\texttt{regression.py} & Main regressions: baselines, tests, in-sample and cross-validated explanatory power. \\
		\texttt{bootstrap.py} & Two-stage bootstrap of both stages, resumable, for any configuration. \\
		\texttt{robustness.py} & Sample and outcome variations, fixed and random effects and the Hausman test. \\
		\texttt{sensitivity\_and\_benchmark.py} & Sensitivity to $\rho$ and $m$ and the principal component benchmark. \\
		\texttt{spec\_grid.py} & Re-runs clustering, alignment and tests for the 120 specifications of the grid. \\
		\texttt{alignment\_check.py} & Checks the consistency of the archetype labels across seasons in all 120 clustering runs. \\
		\texttt{estimated\_weights\_check.py} & Fits the model with estimated attribute weights for one season in four variants (Appendix~\ref{app:weights}). \\
		\texttt{seed\_sensitivity.py} & Repeats the clustering stage with 30 random seeds per fuzziness value and the key test each time. \\
		\texttt{extra\_results.py} & Complete coefficient tables and the baseline without performance statistics. \\
		\texttt{make\_tables\_figures.py} & Produces the tables and figures of the thesis. \\
		\bottomrule
	\end{tabular}
\end{table}

The estimated-weights results of Appendix~\ref{app:weights} are reproduced by \texttt{estimated\_weights\_check.py}; the options that it uses are part of the model module (\texttt{fixed\_weights=None}, the scaling \texttt{norm} set to \texttt{max} or \texttt{mean}, and the noise rule \texttt{dave}). The programs require Python 3 with \texttt{numpy}, \texttt{pandas}, \texttt{scipy}, \texttt{scikit-learn}, \texttt{statsmodels}, \texttt{linearmodels}, \texttt{joblib}, \texttt{matplotlib}, \texttt{soccerdata}, \texttt{requests}, \texttt{beautifulsoup4}, \texttt{rapidfuzz} and \texttt{unidecode}.

\chapter{Estimated attribute weights}
\label{app:weights}

Table~\ref{tab:weights} reports the attribute weights that the estimation of \citeA{durso2022robust} produced for the 2023/24 season ($n=406$ players, $m=1.3$, $\rho=0.5$, ten random starts) in the four variants that were tried: distances scaled by their maximum, as in the original paper, or by their mean, with raw or shrunk success rates. This season was used to develop the model, and the other seasons were not examined in this way. The values were recomputed on the final panel and agree with those of the development runs to the third decimal.

\begin{table}[ht]
	\centering
	\small
	\caption{Estimated attribute weights, season 2023/24}
	\label{tab:weights}
	\begin{tabular}{llrrrrrrr}
		\toprule
		Scaling & Rates & $K$ & Performance & Success rates & Position & Noise & XB & Degenerate \\
		\midrule
		Max & Raw & 3 & 0.038 & 0.959 & 0.003 & 5.2\% & 0.21 & yes \\
		Max & Raw & 4 & 0.031 & 0.967 & 0.002 & 4.7\% & 0.22 & yes \\
		Max & Raw & 5 & 0.000 & 0.000 & 1.000 & 0.0\% & 0.00 & yes \\
		Max & Shrunk & 3 & 0.312 & 0.601 & 0.088 & 34.0\% & 0.34 & no \\
		Max & Shrunk & 4 & 0.238 & 0.704 & 0.059 & 21.4\% & 0.86 & no \\
		Max & Shrunk & 5 & 0.000 & 0.000 & 1.000 & 0.0\% & 0.00 & yes \\
		Mean & Raw & 3 & 0.151 & 0.742 & 0.107 & 6.2\% & 0.42 & no \\
		Mean & Raw & 4 & 0.091 & 0.038 & 0.871 & 16.0\% & 0.08 & yes \\
		Mean & Raw & 5 & 0.000 & 0.000 & 1.000 & 0.0\% & 0.00 & yes \\
		Mean & Shrunk & 3 & 0.128 & 0.109 & 0.762 & 4.4\% & 0.14 & no \\
		Mean & Shrunk & 4 & 0.052 & 0.044 & 0.905 & 0.0\% & 0.05 & yes \\
		Mean & Shrunk & 5 & 0.053 & 0.046 & 0.901 & 0.0\% & 6.53 & yes \\
		\bottomrule
	\end{tabular}
	\par\smallskip
	\parbox{0.95\textwidth}{\footnotesize\emph{Note.} Weights sum to one. Noise is the share of players whose largest membership is in the noise cluster and XB the Xie--Beni index. A solution is called degenerate if one attribute group has a weight above 0.8, if the Xie--Beni index is below $10^{-6}$, or if more than half of the players are in the noise cluster. A weight of one on the positions with a noise share of zero is the case in which the clusters reproduce the position labels exactly. With the maximum, the two success rates dominate; with the mean, positions dominate. The weights used in the thesis are fixed at 0.501, 0.224 and 0.275 (Section~\ref{sec:method-model}). The table is reproduced by \texttt{estimated\_weights\_check.py}.}
\end{table}

\chapter{Coefficient tables}
\label{app:coef}

Tables~\ref{tab:coef-core13} and~\ref{tab:coef-full13} report all coefficients of the extended models for the main configuration, with cluster-robust standard errors (by player) that treat the memberships as observed. Statistics are in the scaled form of Section~\ref{sec:data-variables}, in units of their average absolute deviation. The numbering of the archetypes is inherited from the first season and differs between runs. In the main configuration ($m=1.3$) archetype 1 is the defenders (the reference), 2 the strikers, 3 the wide attackers and 4 the central midfielders; the same holds at $m=1.4$ (Table~\ref{tab:coef-core14}). At $m=1.5$ the defenders are again archetype 1, and archetype 2 is the wide attackers, 3 the central midfielders and 4 the strikers.

\begin{table}[ht]
	\centering
	\footnotesize
	\caption{Core baseline with membership changes, $m=1.3$ (main configuration)}
	\label{tab:coef-core13}
	\begin{tabular}{lrrr}
\toprule
Variable & Coef. & SE & p \\
\midrule
Age & -0.162 & 0.036 & 0.000 \\
Age squared & 0.002 & 0.001 & 0.005 \\
Log minutes (level) & 0.188 & 0.025 & 0.000 \\
Change in log minutes & 0.372 & 0.022 & 0.000 \\
Defender & 0.013 & 0.031 & 0.673 \\
Forward & -0.042 & 0.042 & 0.312 \\
Team points per match (level) & 0.006 & 0.026 & 0.815 \\
Change in team points per match & 0.288 & 0.033 & 0.000 \\
Change of club within the league & -0.012 & 0.036 & 0.728 \\
Non-penalty goals per 90 (level) & 0.069 & 0.014 & 0.000 \\
Non-penalty goals per 90 (change) & 0.061 & 0.013 & 0.000 \\
Assists per 90 (level) & 0.003 & 0.013 & 0.817 \\
Assists per 90 (change) & 0.023 & 0.012 & 0.064 \\
Pair starting 2018/19 & -0.131 & 0.046 & 0.004 \\
Pair starting 2019/20 & 0.302 & 0.053 & 0.000 \\
Pair starting 2020/21 & -0.005 & 0.045 & 0.910 \\
Pair starting 2021/22 & -0.002 & 0.052 & 0.964 \\
Pair starting 2022/23 & 0.145 & 0.050 & 0.004 \\
Pair starting 2023/24 & 0.079 & 0.048 & 0.097 \\
Change in membership of archetype 2 & -0.036 & 0.102 & 0.723 \\
Change in membership of archetype 3 & -0.216 & 0.083 & 0.010 \\
Change in membership of archetype 4 & -0.104 & 0.072 & 0.147 \\
Change in noise membership & -0.141 & 0.132 & 0.285 \\
\bottomrule
\end{tabular}

\end{table}

\begin{table}[ht]
	\centering
	\scriptsize
	\caption{Full baseline with membership changes, $m=1.3$ (main configuration)}
	\label{tab:coef-full13}
	\begin{tabular}{lrrr}
\toprule
Variable & Coef. & SE & p \\
\midrule
Age & -0.159 & 0.036 & 0.000 \\
Age squared & 0.002 & 0.001 & 0.006 \\
Log minutes (level) & 0.193 & 0.027 & 0.000 \\
Change in log minutes & 0.375 & 0.024 & 0.000 \\
Defender & 0.049 & 0.039 & 0.210 \\
Forward & -0.030 & 0.050 & 0.549 \\
Team points per match (level) & 0.006 & 0.027 & 0.818 \\
Change in team points per match & 0.296 & 0.033 & 0.000 \\
Change of club within the league & -0.011 & 0.037 & 0.763 \\
Non-penalty goals per 90 (level) & 0.050 & 0.039 & 0.203 \\
Non-penalty goals per 90 (change) & 0.033 & 0.030 & 0.267 \\
Assists per 90 (level) & 0.002 & 0.015 & 0.883 \\
Assists per 90 (change) & 0.025 & 0.013 & 0.052 \\
Pair starting 2018/19 & -0.122 & 0.047 & 0.009 \\
Pair starting 2019/20 & 0.312 & 0.056 & 0.000 \\
Pair starting 2020/21 & 0.012 & 0.047 & 0.803 \\
Pair starting 2021/22 & 0.008 & 0.056 & 0.890 \\
Pair starting 2022/23 & 0.175 & 0.054 & 0.001 \\
Pair starting 2023/24 & 0.113 & 0.053 & 0.033 \\
Level of shots p90 & -0.015 & 0.057 & 0.795 \\
Change in shots p90 & 0.064 & 0.045 & 0.160 \\
Level of sot p90 & 0.050 & 0.073 & 0.489 \\
Change in sot p90 & -0.024 & 0.050 & 0.624 \\
Level of crosses p90 & -0.003 & 0.011 & 0.766 \\
Change in crosses p90 & 0.010 & 0.016 & 0.520 \\
Level of interceptions p90 & 0.017 & 0.018 & 0.345 \\
Change in interceptions p90 & 0.001 & 0.019 & 0.960 \\
Level of tackles won p90 & 0.007 & 0.015 & 0.608 \\
Change in tackles won p90 & 0.008 & 0.013 & 0.569 \\
Level of fouls p90 & 0.008 & 0.015 & 0.611 \\
Change in fouls p90 & -0.010 & 0.015 & 0.528 \\
Level of fouled p90 & 0.023 & 0.018 & 0.218 \\
Change in fouled p90 & 0.049 & 0.022 & 0.024 \\
Level of offsides p90 & 0.003 & 0.014 & 0.851 \\
Change in offsides p90 & -0.009 & 0.016 & 0.569 \\
Level of sot rate shr & -0.046 & 0.031 & 0.134 \\
Change in sot rate shr & 0.005 & 0.024 & 0.830 \\
Level of g sh shr & 0.022 & 0.030 & 0.458 \\
Change in g sh shr & 0.027 & 0.024 & 0.276 \\
Change in membership of archetype 2 & -0.114 & 0.113 & 0.313 \\
Change in membership of archetype 3 & -0.308 & 0.089 & 0.001 \\
Change in membership of archetype 4 & -0.140 & 0.072 & 0.053 \\
Change in noise membership & -0.239 & 0.161 & 0.138 \\
\bottomrule
\end{tabular}

\end{table}

\begin{table}[ht]
	\centering
	\footnotesize
	\caption{Core baseline with membership changes, $m=1.4$}
	\label{tab:coef-core14}
	\begin{tabular}{lrrr}
\toprule
Variable & Coef. & SE & p \\
\midrule
Age & -0.160 & 0.036 & 0.000 \\
Age squared & 0.002 & 0.001 & 0.005 \\
Log minutes (level) & 0.186 & 0.025 & 0.000 \\
Change in log minutes & 0.377 & 0.022 & 0.000 \\
Defender & 0.013 & 0.031 & 0.676 \\
Forward & -0.043 & 0.042 & 0.298 \\
Team points per match (level) & 0.006 & 0.026 & 0.815 \\
Change in team points per match & 0.286 & 0.033 & 0.000 \\
Change of club within the league & -0.013 & 0.036 & 0.715 \\
Non-penalty goals per 90 (level) & 0.069 & 0.014 & 0.000 \\
Non-penalty goals per 90 (change) & 0.057 & 0.013 & 0.000 \\
Assists per 90 (level) & 0.002 & 0.013 & 0.871 \\
Assists per 90 (change) & 0.017 & 0.012 & 0.151 \\
Pair starting 2018/19 & -0.137 & 0.045 & 0.003 \\
Pair starting 2019/20 & 0.294 & 0.054 & 0.000 \\
Pair starting 2020/21 & 0.006 & 0.046 & 0.896 \\
Pair starting 2021/22 & -0.004 & 0.051 & 0.938 \\
Pair starting 2022/23 & 0.140 & 0.051 & 0.006 \\
Pair starting 2023/24 & 0.076 & 0.048 & 0.113 \\
Change in membership of archetype 2 & -0.061 & 0.105 & 0.563 \\
Change in membership of archetype 3 & -0.258 & 0.088 & 0.003 \\
Change in membership of archetype 4 & -0.123 & 0.076 & 0.105 \\
Change in noise membership & -0.035 & 0.147 & 0.810 \\
\bottomrule
\end{tabular}

\end{table}
